% ReasonControl --- AAAI-style manuscript.
%
% NOTE ON STYLE FILE: this compiles standalone with the `article` class so it
% can be built anywhere. To switch to the AAAI camera-ready format, replace the
% \documentclass line and preamble down to \begin{document} with the official
% aaai26.sty preamble and keep the body unchanged.
%
% Build:  pdflatex main && pdflatex main
\documentclass[letterpaper,10pt,twocolumn]{article}

\usepackage[utf8]{inputenc}
\usepackage[T1]{fontenc}
\usepackage{lmodern}
\usepackage{amsmath}
\usepackage{amssymb}
\usepackage{graphicx}
\usepackage{booktabs}
\usepackage{url}
\usepackage{caption}
\usepackage[margin=0.75in]{geometry}
\setlength{\columnsep}{0.25in}

\title{\bf Reading the Residual Stream to Stop Thinking:\\
Calibrated Convergence Probes for Training-Free\\ Test-Time Compute Control}
\author{Anonymous Submission}
\date{}

\begin{document}
\maketitle

\begin{abstract}
Reasoning language models spend much of their test-time compute after their
answer is already determined, so a sensor that detects convergence could
recover it. We show that convergence is linearly decodable from the residual
stream of DeepSeek-R1-Distill models at 1.5B and 7B (AUC $0.882$ / $0.854$),
that the signal survives a strict position-controlled baseline within every
position stratum (margin $0.147$ / $0.246$), that it is well calibrated
(ECE $0.022$), and that it costs one matrix--vector product on activations the
model has already computed. Used as an early-exit controller under a
pre-registered protocol---dev-only tuning, single-shot test evaluation,
problem-level cluster bootstrap---it removes $52.4\%$ (1.5B) and $58.3\%$ (7B)
of thinking tokens; at 7B it dominates a compute-matched static budget
($27.1\%$ fewer tokens at equal-or-better accuracy), on GPQA-Diamond it
\emph{exceeds} unconstrained reasoning at both scales, and an operating point
selected on the 1.5B transfers to the 7B without retuning.

The accompanying accuracy cost, however, is larger than our pre-registered
non-inferiority margin, and the reason generalizes beyond this system. Once
the probe drives exits rather than merely observing them, its confidence stops
carrying information: realized accuracy is flat at $\approx 0.69$ across every
asserted-confidence bin from $0.75$ to $0.97$, despite excellent on-policy
calibration. A steering actuator built from directions in the same
representation space---individually interpretable under a logit lens---fails
the same way, costing $0.215$ accuracy while increasing token count. We show
that threshold tuning cannot repair either failure, and argue the shared cause
is that both were validated on states produced by undisturbed decoding and
then applied to states produced by their own interventions. Representations
that are well-behaved when observed are not thereby well-behaved when acted
upon; we suggest on-policy probe retraining as the indicated remedy.
\end{abstract}

\section{Introduction}

Reasoning-tuned language models produce long chains of thought whose length is
only loosely coupled to problem difficulty. A substantial fraction of that
computation is spent after the model's answer is, in effect, already fixed:
the trailing tokens re-derive, re-verify, and restate. If a system could tell
\emph{when} the answer has converged, it could stop early and recover that
compute.

Existing approaches to this problem largely infer convergence from behavior:
they decode a trial answer and check agreement
\cite{dynasor,deer}, monitor
specialized neurons \cite{neat}, or train opaque classifiers on
embeddings \cite{aic}. Behavioral probes pay extra decoding; neuron
heuristics and black-box classifiers are hard to audit and offer little
insight into what is being measured.

We take a representational route. We train a linear probe on the residual
stream to predict a well-defined target---whether truncating the chain
\emph{now} and forcing an answer would yield the same answer the model
eventually produces---and use it as a sensor inside a decoding loop. The probe
reads states the forward pass already produced, so the sensor is essentially
free, and because it is linear and calibrated it can be inspected, ablated,
and compared against position-only baselines.

By every standard offline measure, this sensor is good. It separates converged
from unconverged boundaries at AUC $0.882$, it beats a position-only baseline
within every position stratum, it is well calibrated (ECE $0.022$), and the
threshold chosen for one model transfers to a model five times larger without
retuning. Put it in the loop and it removes more than half of all thinking
tokens.

And yet the accuracy cost of those savings is larger than we pre-registered as
acceptable, and the reason is the finding we think matters most here. Once the
probe \emph{drives} exits rather than merely observing them, its confidence
stops carrying information: realized accuracy is flat at $\approx 0.69$ across
every asserted-confidence bin from $0.75$ to $0.97$. The sensor was measured on
states produced by undisturbed decoding, and is then queried on states produced
by its own decisions.

The same thing happens on the write side. We also built a steering actuator
from directions in the same representation space---diff-of-means vectors that
are individually interpretable under a logit lens and behave as intended when
validated in isolation. Applied closed-loop over long horizons it fails
decisively, costing $0.215$ accuracy while \emph{increasing} token count.

Two channels, one pattern, and it organizes this paper:

\begin{quote}
\emph{Representations that are well-behaved when you observe them are not
thereby well-behaved when you act on them.}
\end{quote}

\paragraph{Contributions.}
\begin{enumerate}
\itemsep2pt
\item \textbf{A sensor that is good by every offline measure}: convergence is
linearly and calibratedly decodable beyond position at two scales, at the cost
of one matrix--vector product, with operating points that transfer across
scale (Sec.~\ref{sec:sensor}).
\item \textbf{Acting on it works, and we quantify how well}: $52$--$58\%$ of
thinking tokens removed; at 7B this dominates compute-matched static
budgeting; on GPQA-Diamond it beats unconstrained reasoning outright
(Sec.~\ref{sec:acting}).
\item \textbf{Acting on it also breaks it, in both channels}: a quantitative
demonstration that on-policy calibration is not evidence of calibration under
control, the parallel open-loop/closed-loop failure of the steering channel,
and an argument that threshold tuning cannot repair either
(Sec.~\ref{sec:breaks}). We believe this is a hazard for probe-based
controllers generally, not a property of our system.
\item \textbf{Two measurement pitfalls that silently invert conclusions}
(Sec.~\ref{sec:pitfalls}): answer-extraction coverage and answer-budget
asymmetry between inference engines. Correcting them changed our headline
result twice.
\end{enumerate}

\section{Related Work}

\paragraph{Behavioral early exit.} Dynasor/Certaindex \cite{dynasor}
periodically forces an answer and exits on agreement; DEER \cite{deer}
uses trial-answer confidence; SpecExit \cite{specexit} and FlashThink
\cite{flashthink} use auxiliary models or speculative signals; a related line
stops when successive answers agree \cite{answerconv}. These are effective but
pay additional decoding per check. Our probe consumes activations already
computed.

\paragraph{Latent-state controllers.} NEAT \cite{neat} monitors exit
neurons and suppresses reflection tokens; AIC \cite{aic} trains gradient-boosted
trees on embeddings to choose among continue/terminate/intervene; STU-PID
\cite{stupid} modulates steering strength with a PID controller over a
redundancy classifier; TPV \cite{tpv} probes reasoning progress and steers in
the same direction; CGRS \cite{cgrs} gates reflection suppression on model
certainty at the logit level. Closed-loop, state-conditioned intervention is
therefore not itself novel. Our contribution is a sensor with defined semantics and
audited calibration, an evaluation of every family in a single harness under
one protocol, and an explicit account of when the approach fails. We are
careful about the strength of the baseline comparison: the budget families
received dev sweeps, but our trial-decode implementation did not, and we
qualify that comparison accordingly.

\paragraph{Length control by training.} L1/LCPO \cite{l1}, ShorterBetter, and
O1-Pruner train models to be concise. These are complementary: they change the
model, we change only inference, and the two could compose.

\paragraph{Reasoning-phase representations.} Venhoff et al.\ \cite{venhoff}
characterize cognitive phases (deduction, verification, backtracking) as
linearly represented and steerable. We reuse that taxonomy for a second probe
channel and for a steering actuator, and report a negative causal result for
the actuator.

\section{Method}

\subsection{Chunk boundaries and the control loop}

We decode with a batched manual loop and evaluate the controller only at
\emph{reasoning-chunk boundaries}---double-newline delimited segments of the
chain of thought, which correspond to natural planning units. At each boundary
the controller reads the residual stream at a fixed probe layer, evaluates the
linear probes, and emits one of \textsc{exit}, \textsc{steer}, or
\textsc{continue}. \textsc{exit} splices the end-of-thinking delimiter and
decodes the answer segment greedily; \textsc{continue} resumes sampling.

The probe layer is kept strictly below the steering layer so that the sensor
cannot read its own actuator's edit within the same forward pass, and steering
directions are orthogonalized against the probe weight vectors (post-hoc
cosines $\le 10^{-3}$ in magnitude at both scales).

\subsection{Convergence labels without human annotation}
\label{sec:labels}

The convergence target is defined operationally. For every chunk boundary $t$
in a stored rollout we truncate the chain at $t$, append the end-of-thinking
delimiter, and decode an answer greedily. The boundary is labeled
\emph{converged} if that forced answer matches the answer the unmodified
rollout eventually produced. The label is therefore a property of the model's
own computation, requires no human annotation and no external judge, and
directly matches the decision the controller must make. Parse rate of forced
answers was $1.000$ across all datasets and both models; labels were collected
at $176{,}448$ (1.5B) and $158{,}961$ (7B) boundaries.

The secondary phase channel (exploration / deduction / verification /
backtracking / other) is labeled by regex and audited against a
Qwen2.5-32B-AWQ judge; agreement is moderate (Cohen's $\kappa$ $0.43$--$0.52$),
so phase-derived results are reported as secondary and the primary convergence
labels are judge-independent by construction.

\subsection{Probes}

Probes are logistic regressions on the residual stream at a single layer,
fit on the probe-train split only. Because a naive convergence probe could be
solved by position alone (later boundaries are more likely converged), we
pre-registered a position-controlled gate: the probe must exceed a
position-only shallow baseline by $\ge 0.05$ AUC \emph{within every position
stratum}, and reach AUC $\ge 0.75$ on MATH dev, before any controller
experiment was permitted to run.

\subsection{Policies}

\begin{itemize}
\itemsep2pt
\item \textsc{exit\_only}: \textsc{exit} when the convergence probability
exceeds $\tau$ for $K$ consecutive boundaries (patience), else \textsc{continue}.
\item \textsc{steer\_only}: add $\alpha \hat{v}$ at the steering layer when a
sustained verification/backtracking loop is detected, without exiting.
\item \textsc{full}: both channels.
\item \textsc{noop}: unmodified decoding (the control condition).
\end{itemize}

Baselines implemented in the same harness: \textsc{static\_budget} (hard token
cap), \textsc{budget\_prompt} (budget stated in the prompt),
\textsc{concise\_prompt}, and \textsc{trial\_decode} (a Dynasor-style
forced-answer agreement exit).

\section{Experimental Setup}

\paragraph{Models and data.} DeepSeek-R1-Distill-Qwen-1.5B (probe layer L15,
steer L21) and -7B (L18, L21), fp16. Splits are fixed at seed 0: MATH-train
500 probe-train / 200 dev, MATH-500 as untouched test; GSM8K 300/50/250;
AIME'24 dev and AIME'25 test ($n{=}30$, post-dating the model release);
GPQA-Diamond ($n{=}198$) test-only, with operating points inherited from MATH
dev and no per-dataset tuning. Sampling is temperature $0.6$, top-$p$ $0.95$,
with a 16k thinking cap (24k for AIME) and a 2048-token answer budget
(Sec.~\ref{sec:pitfalls}).

\paragraph{Protocol.} All tuning uses dev data only. The test split is
evaluated once, at the dev-selected operating point, chosen by a 1-SE rule
(best pooled dev accuracy, then fewest tokens among points within one standard
error). Headline conditions use 8 seeds at 1.5B and 4 at 7B.

\paragraph{Statistics.} The unit of analysis is the problem; conditions share
problems and per-(problem, rollout) seeds, so deltas are paired. All intervals
are $95\%$ problem-level cluster bootstraps with $10{,}000$ resamples. The
pre-registered primary endpoint is pooled MATH-500 + GSM8K ($n{=}750$
problems), tested hierarchically: first token superiority, then accuracy
non-inferiority at a $2\%$ margin. AIME'25 and GPQA carry descriptive
intervals only and are never used for tuning.

\section{The Sensor Works}
\label{sec:sensor}

\subsection{Convergence is linearly decodable beyond position}
\label{sec:probes}

Table~\ref{tab:probes} reports probe quality. Both models clear the
pre-registered gate comfortably. AUC peaks in mid-network ($L15/28 \approx
0.54$ depth at 1.5B; $L18/28 \approx 0.64$ at 7B), and---importantly---the
margin over the position-only baseline \emph{grows} with scale ($0.147
\rightarrow 0.246$), which is the opposite of what a position artifact would
predict. On-policy calibration is excellent at 1.5B (ECE $0.022$) and moderate
at 7B ($0.055$). The phase channel reaches macro AUC $0.938$ / $0.937$.

\begin{table}[t]
\centering
\small
\begin{tabular}{lcccc}
\toprule
Model & Layer & AUC & Position margin & ECE \\
\midrule
1.5B & L15 & $0.882$ & $0.147$ & $0.022$ \\
7B   & L18 & $0.854$ & $0.246$ & $0.055$ \\
\bottomrule
\end{tabular}
\caption{Convergence probe quality at the selected layer. ``Position margin''
is the minimum AUC advantage over a position-only baseline across position
strata; the pre-registered gate required $\ge 0.05$.}
\label{tab:probes}
\end{table}

\subsection{Operating points transfer across scale}
\label{sec:transfer}

The 7B results in Table~\ref{tab:primary} use the threshold selected on 1.5B
dev data, with \emph{no} retuning---no 7B sweep was run. It transfers usefully:
the controller remains far cheaper than unconstrained decoding and dominates
static budgeting. It also transfers imperfectly, and informatively: the
accuracy cost roughly doubles ($-0.039 \rightarrow -0.073$) for a comparable
token saving, i.e.\ a threshold calibrated on a weaker model exits slightly too
eagerly for a stronger one. This is consistent with the 7B probe's higher ECE
($0.055$ vs $0.022$) and suggests threshold transfer should be accompanied by
a light recalibration when the target model's probe is less well calibrated.

\section{Acting On It Saves Compute}
\label{sec:acting}

\subsection{Main efficiency result}
\label{sec:main}

Table~\ref{tab:primary} reports the pre-registered primary endpoint. The
controller removes $52.4\%$ of thinking tokens at 1.5B and $58.3\%$ at 7B,
with token superiority significant by a wide margin at both scales. The
accompanying accuracy change is $-0.039$ $[-0.054,-0.024]$ and $-0.073$
$[-0.091,-0.056]$: small in absolute terms, but resolvable at this sample size
and larger than the $2\%$ non-inferiority margin we pre-registered. We
therefore describe the method as trading a small, measurable accuracy cost for
a large compute saving, rather than as accuracy-preserving.

\begin{table}[t]
\centering
\small
\setlength{\tabcolsep}{3pt}
\begin{tabular}{lcc}
\toprule
Pooled MATH-500 + GSM8K & Acc.\ [95\% CI] & Tokens \\
\midrule
\multicolumn{3}{l}{\emph{1.5B}}\\
\textsc{exit\_only} ($\tau{=}0.7,K{=}2$) & $0.817$ $[0.796,0.836]$ & $1{,}773$ \\
\textsc{noop} (same engine) & $0.856$ $[0.835,0.876]$ & $3{,}722$ \\
\textsc{static\_budget} & $0.803$ $[0.780,0.825]$ & $1{,}578$ \\
\midrule
\multicolumn{3}{l}{\emph{7B (transferred point)}}\\
\textsc{exit\_only} & $0.863$ $[0.843,0.882]$ & $1{,}093$ \\
\textsc{noop} (same engine) & $0.936$ $[0.917,0.953]$ & $2{,}622$ \\
\textsc{static\_budget} & $0.849$ $[0.826,0.870]$ & $1{,}499$ \\
\bottomrule
\end{tabular}
\caption{Pre-registered primary endpoint. Paired token reduction is
$-1{,}949$ $[-2{,}106,-1{,}798]$ at 1.5B and $-1{,}529$ $[-1{,}699,-1{,}368]$
at 7B; paired accuracy change is $-0.039$ $[-0.054,-0.024]$ and $-0.073$
$[-0.091,-0.056]$.}
\label{tab:primary}
\end{table}

\paragraph{Against a compute-matched baseline.} The comparison that matters
for a practitioner is not against unconstrained decoding but against simply
capping the budget. Here the result is scale-dependent and, at 7B, clean:
the controller uses $27.1\%$ fewer tokens than \textsc{static\_budget}
($-406$ $[-499,-301]$) at equal-or-better accuracy ($+0.014$
$[+0.000,+0.029]$)---it dominates on both axes. At 1.5B it does not dominate:
it buys $+0.014$ $[+0.001,+0.028]$ accuracy while spending $12\%$ more tokens.
The natural reading is that a learned exit rule pays off once the model is
strong enough that the right stopping point genuinely varies across problems;
on a weaker model, where most chains run to a similar length, a fixed cap
captures most of the available benefit.

\subsection{Per-dataset behavior}

Table~\ref{tab:perds} gives the test grid at 1.5B; the same ordering holds at
7B. The accuracy cost tracks how much genuine search a problem requires:
negligible on GSM8K, moderate on MATH-500, largest on AIME. The prompt-level
baselines are included for reference rather than as competitors---
\textsc{budget\_prompt} barely restrains generation at all ($4{,}514$ tokens
on MATH-500 against \textsc{noop}'s $4{,}653$), so it preserves accuracy
without being an efficiency method.

We exclude our \textsc{trial\_decode} implementation from the comparison. Its
exit rule has free parameters (agreement count, minimum chunks before probing)
that we ran at a single default configuration rather than sweeping, whereas
\textsc{exit\_only} received nine dev configurations and the budget families
four each; the resulting operating point is an order of magnitude more
aggressive than any we selected for ourselves ($317$ vs $2{,}312$ mean think
tokens). Our numbers for it therefore fall far below the near-lossless results
reported by DEER \cite{deer} and Dynasor \cite{dynasor} and should not be read
as an evaluation of that family. A properly tuned head-to-head against
behavioral exit remains open work.

\begin{table}[t]
\centering
\small
\setlength{\tabcolsep}{4pt}
\begin{tabular}{lcccc}
\toprule
1.5B, test & MATH-500 & GSM8K & GPQA-D & AIME'25 \\
\midrule
\textsc{noop} & $.833$ & $.877$ & $.240$ & $.225$ \\
 & \tiny$4653$ & \tiny$2082$ & \tiny$8027$ & \tiny$14604$ \\
\textsc{exit\_only} & $.807$ & $.836$ & $\mathbf{.323}$ & $.237$ \\
 & \tiny$2312$ & \tiny$694$ & \tiny$6538$ & \tiny$11492$ \\
\textsc{static\_budget} & $.774$ & $.861$ & $.316$ & $.146$ \\
 & \tiny$1726$ & \tiny$1281$ & \tiny$1923$ & \tiny$2033$ \\
\textsc{budget\_prompt} & $.832$ & $.867$ & $.253$ & $.258$ \\
 & \tiny$4514$ & \tiny$1916$ & \tiny$8015$ & \tiny$14069$ \\
\textsc{concise\_prompt} & $.832$ & $.765$ & $.296$ & $.254$ \\
 & \tiny$3447$ & \tiny$344$ & \tiny$6693$ & \tiny$14090$ \\
\bottomrule
\end{tabular}
\caption{1.5B test results, 8 seeds. Upper number is accuracy, lower (small)
is mean think tokens. AIME'25 ($n{=}30$) is a descriptive anchor with wide
intervals and supports no inferential claim; GPQA ($n{=}198$) is discussed in
Sec.~\ref{sec:gpqa}.}
\label{tab:perds}
\end{table}

\subsection{Early exit improves hard OOD multiple choice}
\label{sec:gpqa}

On GPQA-Diamond---held out entirely from tuning, and a different task format
(multiple choice, natural science)---the controller \emph{exceeds}
unconstrained reasoning at both scales. Comparing within engine, accuracy is
$0.323$ vs $0.249$ at 1.5B ($+0.074$) and $0.433$ vs $0.384$ at 7B ($+0.049$),
while using $16\%$ and $60\%$ fewer thinking tokens respectively. \textsc{static\_budget} also improves over
\textsc{noop} here ($0.316$, $0.405$), so the honest statement is that
restraint of any kind helps on this benchmark rather than that our controller
is uniquely responsible. The finding itself is the point: on hard multiple
choice, long unconstrained chains do not merely waste compute, they degrade
answers---the model reasons itself away from an initially correct choice. A
convergence-aware exit captures that benefit while remaining adaptive.

\section{But Acting Breaks the Sensor}
\label{sec:breaks}

\subsection{Calibration collapses under intervention}
\label{sec:audit}

This is, we believe, the most transferable finding in the paper. The
convergence probe is well calibrated on stored rollouts (ECE $0.022$). Once it
\emph{drives} the policy, that calibration disappears. Table~\ref{tab:audit}
bins exited rollouts by the confidence the probe asserted at the moment of
exit, and reports realized accuracy in each bin. Realized accuracy is
essentially flat at $\approx 0.69$ across the whole range from $0.75$ to
$0.97$: the probe's confidence, once it is in the loop, carries almost no
information about whether the answer is actually right.

\begin{table}[t]
\centering
\small
\begin{tabular}{lccc}
\toprule
Claimed $p$ bin & $n$ & Mean claimed $p$ & Realized acc. \\
\midrule
$[0.70,0.80)$ & $230$ & $0.755$ & $0.670$ \\
$[0.80,0.90)$ & $281$ & $0.852$ & $0.712$ \\
$[0.90,0.95)$ & $206$ & $0.925$ & $0.694$ \\
$[0.95,1.00]$ & $233$ & $0.971$ & $0.691$ \\
\bottomrule
\end{tabular}
\caption{Closed-loop calibration audit, 1.5B at the selected operating point.
The probe is well calibrated on-policy (ECE $0.022$) but its asserted
confidence is uninformative once it drives exits.}
\label{tab:audit}
\end{table}

The mechanism is distribution shift of a specific and avoidable kind. The
probe is trained on states produced by \emph{unintervened} decoding, and then
queried on states produced by its own interventions---a control problem
structurally identical to the one imitation learning addresses with iterative
on-policy data collection. This bounds what threshold tuning can achieve: our
dev sweep over $\tau \in \{0.7,0.8,0.9,0.95\}$ and $K \in \{1,2\}$ is
essentially flat in accuracy, because raising the bar buys tokens rather than
correctness when the bar itself is uninformative. The indicated fix is not a
better threshold but a probe retrained on the states it induces. We regard
this as a caution for the broader class of probe-driven controllers: reporting
on-policy AUC and ECE, as is standard, does not establish that a sensor will
behave when it is put in charge.

\subsection{The write channel fails the same way}
\label{sec:steer}

The read channel is not a special case. We built a steering actuator from
diff-of-means directions for the reasoning phases, orthogonalized against the
probe weights so that sensor and actuator are decoupled. Validated in
isolation these directions are exactly what they claim to be: under a logit
lens the backtracking direction most promotes \emph{but}, \emph{However},
\emph{But}, \emph{perhaps}---including the corresponding Chinese discourse
connectives, indicating a language-general feature rather than a token
pattern---and the verification direction promotes \emph{checking},
\emph{check}, \emph{verification}, \emph{confirming}.

Applied closed-loop, they fail. We pre-registered an acceptance gate---the
actuator ships only if the upper confidence bound on the paired accuracy drop
is below $2\%$ over at least 200 paired dev rollouts---and it fails by an
order of magnitude. At $\alpha{=}3$ over 250 paired problems (1000 rollouts),
steering costs $-0.215$ $[-0.251,-0.180]$ accuracy and, far from saving
compute, \emph{adds} $+3{,}583$ $[+3{,}131,+4{,}063]$ thinking tokens. A
regex screen shows the intervention moving in the wrong direction entirely:
positive $\alpha$ \emph{lowered} the backtracking rate ($0.163$ vs $0.231$,
$\Delta = -0.068$ $[-0.074,-0.062]$) rather than suppressing it as intended.
A second pre-registered gate (steering must contribute $\ge 5\%$ of
exit-only's savings) therefore resolved this paper to be exit-led, in writing,
before test results were examined.

Interpretability of a direction, then, is not evidence that the direction is
safe to act on---the same gap that separates the probe's offline calibration
from its behavior in the loop. Notably, the convergence direction is the
\emph{less} interpretable of the two under a logit lens (apart from the
end-of-sequence token, its top promotions are not semantically coherent) yet
it is the one that produces a useful controller. Interpretability and
actionability are close to orthogonal here.

\subsection{Why threshold tuning cannot repair this}
\label{sec:whynot}

A natural response is that the exit threshold is simply set too aggressively.
The dev sweep says otherwise: accuracy is close to flat in $\tau$ across
$\{0.7,0.8,0.9,0.95\}$ and $K \in \{1,2\}$, so raising the bar buys tokens
rather than correctness. This is what Table~\ref{tab:audit} predicts---if
asserted confidence does not track correctness under intervention, then
thresholding that confidence more strictly cannot buy correctness.

We can make the argument quantitative. Estimating each grid point's accuracy
at the corrected answer budget from its unparsed-answer rate, calibrated on
the two configurations we measured at both budgets, the 1-SE selection rule
returns the \emph{same} operating point under both a conservative and a
generous calibration; the best attainable accuracy anywhere on the grid
remains below the unintervened control. The binding constraint is the sensor's
behavior under its own intervention, not the threshold applied to it.

The implied fix is standard in imitation learning and absent from probe-based
control: collect probe training data \emph{on the states the controller
induces} and retrain iteratively, rather than on undisturbed rollouts. We did
not have the budget to do this here, and we think it is the most promising
next step for this class of method.

\section{Two Measurement Pitfalls}
\label{sec:pitfalls}

We report these because both are invisible in aggregate metrics, both inverted
our conclusions, and neither is specific to our system.

\paragraph{Answer-extraction coverage.} Our initial extractor accepted
\verb|\boxed{}| and ``the answer is $X$'' but not the
``\textbf{Answer:} $X$'' form, which these models emit in roughly 20--25\% of
mathematical completions and about 60\% on GPQA. Those rollouts were scored
incorrect. The defect was caught only because two GPQA numbers fell below the
$0.25$ chance floor of a four-way multiple choice---an impossibility that a
purely relative comparison would have hidden. All $26{,}926$ stored rollouts
were re-graded offline from stored text, with no regeneration.

\paragraph{Answer-budget asymmetry across engines.} Our controller ran in a
HuggingFace decode loop that enforced a 512-token answer cap; the vLLM
baselines received thinking and answer budgets jointly and were effectively
uncapped. This truncated exactly the verbose conditions---including the
\textsc{noop} control---and manifested as an apparent $0.16$ ``cross-engine
accuracy gap'' that we initially, and wrongly, attributed to the inference
engine. Raising the answer budget to 2048 tokens (above the observed p99
answer length) eliminated it: the two engines now agree to within $0.013$ on
every dataset at both scales. Any comparison that mixes inference engines, or
that caps generation per-segment rather than per-sequence, should verify this
explicitly.

\section{Limitations}

The operating point was selected on dev data collected under the 512-token
answer budget. We verified the selected point at the corrected budget and
estimated the full grid's behavior there; the 1-SE rule re-selects the same
point under both conservative and generous assumptions, so we do not believe
selection is a confound, but we did not re-run the entire sweep. The 7B
within-engine control uses a single seed (the 7B cross-engine control uses
four). The \textsc{trial\_decode} baseline was run at one default
configuration rather than swept, so our numbers understate that family and we
draw no comparative conclusion from them. AIME'25 ($n{=}30$) and GPQA ($n{=}198$) are descriptive anchors and
carry no inferential claims. Phase labels have moderate judge agreement
($\kappa$ $0.43$--$0.52$), so phase results are secondary throughout. All
results are on the DeepSeek-R1-Distill-Qwen family; transfer to other reasoning
model families is untested.

\section{Conclusion}

Whether a reasoning model has finished thinking is written legibly in its
residual stream. It is readable with a single matrix--vector product on
activations already computed, it survives a strict position control at two
scales, and the threshold that reads it transfers from a 1.5B model to a 7B
one untouched. Acting on it removes more than half of all thinking tokens,
dominates compute-matched static budgeting at 7B, and improves accuracy
outright on hard out-of-distribution multiple choice.

The limits are the more transferable half of the result. The same probe whose
calibration is excellent on stored rollouts asserts confidences that carry
almost no information once it is the thing choosing when to stop; a steering
actuator built from demonstrably interpretable directions in the same space
degrades the chain it was meant to improve. Neither failure is a tuning
problem, and both share a cause: each component was measured on states the
undisturbed model produces and then deployed on states it produces itself.

For work that puts probes in control loops, we think the practical implication
is concrete. Offline AUC, ECE, and interpretability---the standard evidence
that a direction means what we say it means---establish that a representation
is readable, not that it is actionable. Those are different properties, they
can come apart sharply, and closing the gap likely requires training sensors
on the distribution their own decisions create.

\begin{thebibliography}{99}
\small

\bibitem{dynasor}
Fu, Y.; et al.
\newblock Efficiently Serving LLM Reasoning Programs with Certaindex.
\newblock arXiv:2412.20993, 2024.

\bibitem{deer}
Yang, C.; et al.
\newblock Dynamic Early Exit in Reasoning Models.
\newblock arXiv:2504.15895, 2025.

\bibitem{neat}
NEAT: Neuron-level Early-exit And Termination for reasoning models.
\newblock arXiv:2602.02010, 2026.

\bibitem{aic}
Adaptive Inference-time Control for Reasoning Models.
\newblock NeurIPS Workshop, 2025.

\bibitem{stupid}
STU-PID: Steering Token Usage via PID Control.
\newblock arXiv:2506.18831, 2025.

\bibitem{tpv}
Overclocking LLM Reasoning: Thinking Progress Vectors.
\newblock arXiv:2506.07240, 2025.

\bibitem{specexit}
SpecExit: Speculative Early Exit for Reasoning Models.
\newblock arXiv:2509.24248, 2025.

\bibitem{flashthink}
FlashThink: An Early Exit Method for Efficient Reasoning.
\newblock arXiv:2505.13949, 2025.

\bibitem{l1}
Aggarwal, P.; and Welleck, S.
\newblock L1: Controlling How Long a Reasoning Model Thinks with
Reinforcement Learning.
\newblock arXiv:2503.04697, 2025.

\bibitem{venhoff}
Venhoff, C.; et al.
\newblock Understanding Reasoning in Thinking Language Models via Steering
Vectors.
\newblock arXiv:2506.18167, 2025.

\bibitem{cgrs}
Certainty-Gated Reflection Suppression.
\newblock arXiv:2508.05337, 2025.

\bibitem{answerconv}
Answer Convergence as a Signal for Early Stopping.
\newblock arXiv:2506.02536, 2025.

\end{thebibliography}

\end{document}
